% !PRINTED 359
\weberpage{359}
Fallen die Doppeltangenten in eine Linie zusammen, so treten zwei Doppelpunkte auf.

Dagegen können sehr wohl, ohne dass singuläre Punkte entstehen, die beiden Berührungspunkte einer Doppeltangente zusammenfallen, und so die Doppeltangente in eine vierpunktig berührende Tangente übergehen. Dies tritt ein, wenn die Linie $x_1$ den Kegelschnitt $u$ berührt.

Auch die Form (2) ist noch mit Beibehaltung von $x_1, y_1$ auf unendlich viele Arten herzustellen.

Nehmen wir, um dies einzusehen, an, es sei
\[
f = x_1 y_1 v - u^2 = x_1 y_1 v_1 - u_1^2,
\]
so folgt die identische Gleichung:
\weq{3}{x_1 y_1 (v - v_1) = (u - u_1)(u + u_1).}
Wenn nun $u - u_1$ durch $x_1$, $u + u_1$ durch $y_1$ theilbar wäre, so würde im Schnittpunkte von $x_1$ und $y_1$ auch $u$ und folglich $f$ verschwinden. Dieser Schnittpunkt würde auf der Curve $f$ liegen, was, so lange die Curve keinen singulären Punkt hat, nicht möglich ist. Es muss also einer der beiden Factoren $u - u_1$, $u + u_1$ in (3) durch $x_1 y_1$ theilbar sein. Da das Vorzeichen von $u_1$ noch nicht bestimmt ist, so können wir annehmen, es sei $u - u_1$ durch $x_1 y_1$ theilbar, und also bis auf einen constanten Factor mit $x_1 y_1$ identisch. Es wird dann, wenn $\lambda$ ein solcher Factor ist,
\[
u_1 = u + \lambda x_1 y_1,
\]
und die Function $f$ erhält die Form:
\weq{4}{f = x_1 y_1 (v + 2\lambda u + \lambda^2 x_1 y_1) - (u + \lambda x_1 y_1)^2.}
Umgekehrt sind für jedes beliebige $\lambda$ die Formeln (2) und (4) mit einander identisch.

Wenn wir nun $\lambda$ so bestimmen können, dass die quadratische Function $v + 2\lambda u + \lambda^2 x_1 y_1$ in zwei lineare Factoren zerfällt, so erhält, wenn wir $u + \lambda x_1 y_1 = u_1$ setzen, $f$ nach (4) die Gestalt
\weq{5}{f = x_1 y_1 x_2 y_2 - u_1^2,}
und diese Form zeigt, dass $x_2, y_2$ \gesperrt{zwei weitere Doppeltangenten sind}, und dass die acht Berührungspunkte der Doppeltangenten $x_1, y_1, x_2, y_2$ auf einem Kegelschnitte $u_1 = 0$ liegen.
